% ===========================================================================
% Appendix C. Network design algorithms
% Specification: R/indices/allied_network_design.R and scripts/30_build_allied_network_design.R
% (production settings). Lambda study: final calibration of 2026-10-12 on the 37-node set
% (scratch_final/network/lambda_study_g37.csv, lambda_study_nc37.csv; paper/results/
% network_frontier.json). Notation as in Section 10.
% ===========================================================================
\section{Network design algorithms}
\label{app:algorithm}

This appendix states the algorithms behind the allied network design of Section~\ref{sec:network},
proves the properties used there, and documents the calibration of the diversification weight.
Notation follows Section~\ref{sec:network}: $\mathcal{V}$ is the set of 37 network nodes,
$a_p$ the producer score, $d_c$ the demand weight, $W_{pc}$ the edge weight, $\delta_p$ the
development potential, $\mathcal{X}_{ts}$ the nodes barred from stage $(t,s)$ by the portfolio caps
and $\mathcal{E}_{ts}$ the eligible producers. The functions of
\texttt{R/indices/allied\_network\_design.R} read and write no files; all input and output is in
\texttt{scripts/30\_build\_allied\_network\_design.R}. Table~\ref{tab:alg_params} lists the
production settings.

\begin{table}[htbp]
\centering
\begin{threeparttable}
\caption{Parameters of the allied network design (production configuration)}
\label{tab:alg_params}
\footnotesize
\begin{tabularx}{\textwidth}{@{}l l l X@{}}
\toprule
Symbol & Code argument & Value & Role \\
\midrule
\multicolumn{4}{@{}l}{\emph{Nodes, demand and edges}}\\
-- & \texttt{network\_members} & 57 + 6, 37 nodes & Coalition members and EU-only members; the EU-27 as node \texttt{EUU} \\
-- & \texttt{supply\_chains} & Up, Mid & The 22 stages solved \\
$\gamma$ & -- & 0.5 & S-curve exponent of node components \\
$\omega_{\EO},\omega_{\ES}$ & \texttt{node\_weights} & 0.5, 0.2 & Producer score, eq.~\eqref{eq:net_producer_score} \\
-- & \texttt{node\_policy} & \texttt{FALSE} & Policy term off \\
-- & \texttt{demand\_mode} & \texttt{mixed} (0.5, 0.5) & Need and size, eq.~\eqref{eq:net_demand} \\
$\beta_{\FS},\beta_{\PO}$ & \texttt{edge\_weights} & 0.5, 0.5 & Edge weight, eq.~\eqref{eq:net_edge} \\
\multicolumn{4}{@{}l}{\emph{Stage program}}\\
$w_N,w_E,w_D$ & \texttt{w\_node}, \texttt{w\_edge}, \texttt{w\_dev} & 1, 0.5, 0 & Objective weights, eq.~\eqref{eq:net_obj} \\
$\lambda$ & \texttt{diversity\_lambda} & 0.78 & Diversification weight (calibrated) \\
$h$ & \texttt{diversity\_breakpoint\_step} & 0.01 & Share grid of the tangent cuts \\
$\bar s,\underline s$ & \texttt{max\_share}, \texttt{min\_share} & 0.40, 0.05 & Share bounds \\
$m_P,m_S$ & \texttt{min\_producers}, \texttt{min\_suppliers\_per\_consumer} & 3, 2 & Producers per stage, suppliers per consumer \\
$\varepsilon$ & \texttt{epsilon\_supplier\_share} & 0.10 & Smallest share of a consumer's demand on an active link \\
\multicolumn{4}{@{}l}{\emph{Solver and portfolio caps}}\\
-- & \texttt{auto\_milp\_max\_nodes} & 60 & Largest stage sent to the MILP \\
$T$ & \texttt{milp\_stage\_time\_limit\_sec} & 300 & Time limit per stage solve (s) \\
$K$ & \texttt{portfolio\_top\_k} & 5 & Top-$K$ rule, eq.~\eqref{eq:net_topk} \\
$\kappa_{\min},\kappa_{\max}$ & \texttt{portfolio\_min\_cap}, \texttt{portfolio\_max\_cap} & 2, 10 & Caps, eq.~\eqref{eq:net_cap}, log-GDP scale \\
$I_{\max}$ & \texttt{portfolio\_max\_iters} & 30 & Largest number of cap passes \\
\bottomrule
\end{tabularx}
\begin{tablenotes}[flushleft]\footnotesize
\item \emph{Notes:} Values set in \texttt{scripts/30\_build\_allied\_network\_design.R} or, where the
script passes none, the defaults of \texttt{allied\_network\_design()}. Component weights are
renormalised over the components observed; objective weights are used as given.
\end{tablenotes}
\end{threeparttable}
\end{table}

% ---------------------------------------------------------------------------
\subsection{Overview and network nodes}
\label{subsec:alg_overview}

Algorithm~\ref{alg:network_design} gives the top level. The stage programs are solved one at a time
and independently; the only link between stages is the portfolio loop, which changes the exclusion
sets and re-solves the stages it affects.

\begin{algorithm}[htbp]
\caption{Allied network design (top level)}
\label{alg:network_design}
\begin{algorithmic}[1]
\Require cell indices $\EO,\ES$; directed partner indices $\FS,\PO$; development table; GDP; members
  and their nodes; EU stage weights; 2024 basket trade among members; parameters of
  Table~\ref{tab:alg_params}
\State collapse every input to one record per key; \textbf{assert} uniqueness
\State aggregate the members of the EU node (rules below); every other member is its own node
\State for each stage, S-curve $\EO$ and $\ES$ across the nodes; compute $a_p$, $d_c$, $W_{pc}$
\State $\kappa_i\gets$ eq.~\eqref{eq:net_cap} for every node; \ $\mathcal{X}_{ts}\gets\emptyset$
\ForAll{stages $(t,s)$}
  \State $R_{ts}\gets$ \textit{SolveStage}$(t,s,\mathcal{X}_{ts})$\Comment{Algorithm~\ref{alg:stage}}
\EndFor
\State $(R,\mathcal{X})\gets$ \textit{PortfolioCaps}$(R,\mathcal{X},\kappa)$\Comment{Algorithm~\ref{alg:portfolio}}
\State split each EU share to the members by stage weight; rank build candidates by
  eq.~\eqref{eq:net_build}
\State \Return designs, diagnostics, caps and counts, member shares, build list, input checks
\end{algorithmic}
\end{algorithm}

\paragraph{The EU node.} The design is solved over nodes, and the members of the EU node are
aggregated before anything is ranked, so that the node is one point in each within-stage S-curve.
The member weights $\mu_{m,ts}$ are the stage weights of Section~\ref{sec:coalition} (members'
production shares for Coal, Gas and Oil Upstream, their basket-export shares where the stage has a
basket, their GDP shares otherwise), renormalised over the members. A producer input of the node
($\EO$, $\ES$, development potential) is the $\mu$-weighted mean over the members that have a value,
the weights renormalised over them. Demand uses the members' summed GDP for size, and for need the
GDP-weighted mean of the members' raw security, which is then S-curved across nodes like any other
node's. An edge component ($\FS$ or $\PO$) between the node and a partner is the mean of the member
edges, weighted by the members' 2024 basket trade on that edge (the members' imports from the
producer when the EU consumes, their exports to the consumer when it produces) or, without trade, by
the stage weights; components are averaged before the missing-component zero is applied. Edges
between two members are dropped, and the node's self-supply edge follows eq.~\eqref{eq:net_self_edge}.
The portfolio cap uses the summed GDP. The node's designed share is split back to the members in
proportion to $\mu_{m,ts}$; these member shares are derived, not designed.

% ---------------------------------------------------------------------------
\subsection{The stage MILP with the diversification term}
\label{subsec:alg_milp}

The quadratic penalty of eq.~\eqref{eq:net_obj} is carried into the mixed-integer linear program by
an outer approximation. Let $\mathcal{B}=\{0,h,2h,\dots,\bar s\}$ be the breakpoints of the share
grid ($h=0.01$, 41 points). Each producer gets an auxiliary variable $\tau_p\ge0$ and the cuts
\begin{equation}\label{eq:alg_tangent}
  \tau_p\ \ge\ 2b\,x_p-b^{2}\,y_p\qquad \forall b\in\mathcal{B},
\end{equation}
and the objective subtracts $\lambda\sum_p\tau_p$ instead of $\lambda\sum_px_p^2$. With $y_p=1$
the cuts are the tangents of $x^2$ at the breakpoints; with $y_p=0$, so that $x_p=0$, they reduce
to $\tau_p\ge0$. Because the program maximises, $\tau_p$ settles on the upper envelope of the
tangents.

\paragraph{Approximation error.} For $x\in[b,b+h]$ the envelope is $\max\{2bx-b^2,\,2(b+h)x-(b+h)^2\}$,
and $x^2$ minus the envelope is largest where the two tangents meet, at $x=b+h/2$, where it equals
$(h/2)^2$. The envelope therefore under-states $x_p^2$ by at most $h^2/4=2.5\times10^{-5}$ per
producer, and by at most $5\times10^{-4}$ of the Herfindahl index even with the largest feasible
number of producers (20); unselected nodes contribute no error. The diagnostics record the realised error; in the production design it equals
the bound in every stage, because an interior optimal share sits on a kink of the envelope, at a
midpoint $b+h/2$ of the grid, unless a bound or a linking constraint binds. Designed shares such as
0.325, 0.275 or 0.135 are of this kind, which is why ties are frequent and the top-$K$ rule uses a
strict rank.

\paragraph{Perspective form.} Written with $y_p$ multiplying $b^2$, the cuts are the perspective of
the tangents. They leave the integer solutions and their objective values unchanged but tighten the
linear relaxation: with plain tangents $\tau_p\ge2bx_p-b^2$, a fractional $y_p$ could hold a share
below $\underline s$ at almost no penalty. On a Nuclear Upstream stage with near-tied candidates
the perspective form cut GLPK's solution time from 875 to 51 seconds with the same solution.

\paragraph{Size.} With $n$ nodes the program has $n^2$ flows, $n$ selection and $n^2$ link binaries,
and $2n$ share and auxiliary variables: $2n^2+3n$ variables, $n^2+n$ of them binary. It has $n$
demand constraints, $2n$ share bounds, one producer-count constraint, $3n^2$ linking constraints,
$n$ redundancy constraints, $n$ share definitions and $|\mathcal{B}|\,n$ tangent cuts, plus one
equality per barred node. For $n=37$ this is 2{,}849 variables (1{,}406 binary) and 5{,}810
constraints.

\paragraph{Solution and fallback.} Algorithm~\ref{alg:stage} solves one stage. A stage is sent to
the MILP when the solver packages are installed and the stage has at most 60 nodes, which every
stage satisfies. GLPK receives its own time limit of $T=300$ seconds, an R-level limit guards the
model build, and only a solution with optimal status is accepted; a time-limited or infeasible run
falls back to the greedy heuristic, and the diagnostics record the method and the solver's message.

\begin{algorithm}[htbp]
\caption{Solving one stage (\textit{SolveStage})}
\label{alg:stage}
\begin{algorithmic}[1]
\Require stage $(t,s)$; exclusion set $\mathcal{X}_{ts}$; node and edge data of the stage
\State $C\gets$ the nodes of the stage; \textbf{if} $|C|<3$ or no $a_p$ is defined \textbf{then
  return} skipped
\State complete $W$: a missing component is $0$; $W_{pp}\gets$ eq.~\eqref{eq:net_self_edge}
\State $\mathcal{E}\gets C\setminus\mathcal{X}_{ts}$, restricted to nodes with a defined $a_p$ if at
  least $\max\{m_P,\lceil1/\bar s\rceil,m_S\}$ remain
\If{the MILP packages are installed \textbf{and} $|C|\le60$}
  \State build \eqref{eq:net_milp} with the cuts \eqref{eq:alg_tangent} and $y_p=0$ for
    $p\notin\mathcal{E}$; solve with GLPK under the limit $T$
  \If{the build or solve fails, or the status is not optimal}
    \State record the message; $R\gets$ \textit{Greedy}$(C,\mathcal{E})$\Comment{Algorithm~\ref{alg:greedy}}
  \Else
    \State $R\gets$ the optimal $(y,f)$ with $x_p=\sum_cf_{pc}$
  \EndIf
\Else
  \State $R\gets$ \textit{Greedy}$(C,\mathcal{E})$
\EndIf
\State record the method, producers, $J$ with the exact penalty $\lambda\sum_px_p^2$, $\HHI$,
  $\max_px_p$, realised redundancy and the tangent error
\State \Return $R$
\end{algorithmic}
\end{algorithm}

% ---------------------------------------------------------------------------
\subsection{Properties of the stage program}
\label{subsec:alg_properties}

Throughout, self-supply is allowed, $d_c\ge0$ and $\sum_cd_c=1$.

\paragraph{Feasibility.} At the production values the program is feasible if and only if
$|\mathcal{E}_{ts}|\ge3$. \emph{Proof.} Summing the demand constraints gives $\sum_px_p=1$, and
$x_p\le\bar s$ requires at least $\lceil1/\bar s\rceil=3$ selected nodes, all eligible. Conversely,
select any three eligible nodes, link each to every consumer and set $f_{pc}=d_c/3$: shares of $1/3$
lie in $[\underline s,\bar s]$, each link carries $d_c/3\ge\varepsilon d_c$, and every consumer has
three suppliers. The penalty does not affect feasibility. $\square$

\paragraph{Concentration bounds.} Any feasible design satisfies eq.~\eqref{eq:net_hhi_bounds}.
\emph{Proof.} With $k$ producers, each with $x_p\ge\underline s$, $k\le\lfloor1/\underline s\rfloor$
and $k\le|\mathcal{E}|$, and $\sum_px_p^2\ge1/k$ by the Cauchy--Schwarz inequality. The Herfindahl
index is convex, so its maximum over $\{x:\sum_px_p=1,\,0\le x_p\le\bar s\}$ is attained at a vertex,
where $\lfloor1/\bar s\rfloor$ shares equal $\bar s$ and one equals the remainder, here
$(0.4,0.4,0.2)$ and 0.36. $\square$

\paragraph{Redundancy.} Every consumer with $d_c>0$ has between $m_S=2$ and
$\lfloor1/\varepsilon\rfloor=10$ suppliers, and at least $\varepsilon d_c$ of its demand comes from
another node, because an active link carries at least $\varepsilon d_c$ and at least two links are
active.

\paragraph{Interior shares.} Fix the producer set $S$ and write $v_p$ for the marginal value of a
unit of $p$'s share (its node score plus the edge value of the marginal flow). Maximising
$\sum_{p\in S}v_px_p-\lambda\sum_{p\in S}x_p^2$ subject to $\sum_px_p=1$ and
$\underline s\le x_p\le\bar s$ gives, by the Karush--Kuhn--Tucker conditions,
\begin{equation}\label{eq:alg_waterfill}
  x_p=\min\Big\{\bar s,\ \max\Big\{\underline s,\ \frac{v_p-\mu}{2\lambda}\Big\}\Big\},
\end{equation}
with the level $\mu$ set so that the shares sum to one. Two interior shares therefore differ by
$(v_p-v_q)/(2\lambda)$, and as $\lambda\to0$ the shares move to the bounds, the corner solutions of
the linear design. Equation~\eqref{eq:alg_waterfill} is the water-filling rule of the greedy
fallback.

\paragraph{Monotonicity.} Holding the eligible set, scores, demand and edges fixed, the optimal value
is weakly decreasing in the stringency of each resilience parameter (lower $\bar s$, higher
$\underline s$, $m_P$, $m_S$ or $\varepsilon$), because each parameter enters one family of
constraints and a more stringent value only removes feasible points. It is also weakly decreasing in
$\lambda$. The portfolio loop can produce different exclusion sets at different parameter values, so
the property holds stage by stage for given exclusions, not for the design as a whole; the
non-monotone median Herfindahl index of the calibration grid arises in this way.

% ---------------------------------------------------------------------------
\subsection{Greedy fallback}
\label{subsec:alg_greedy}

The heuristic (Algorithm~\ref{alg:greedy}) replaces the joint choice of producers, shares and flows
by three steps. It ranks the eligible nodes on $g_p=w_Na_p+w_E\bar W_p+w_D\delta_p$, where
$\bar W_p=\sum_cW_{pc}d_c$ is the demand-weighted connectivity, with ties broken by $a_p$, then
$\bar W_p$, then the code, so that the selection never depends on row order; nodes with no positive
edge to a consumer with demand go last. For each $k\ge m_P$ it water-fills the first $k$ nodes with
eq.~\eqref{eq:alg_waterfill}, using $v_p=g_p$ (the edge term as if $p$ served every consumer in
proportion to demand) and finding $\mu$ by bisection, and keeps the $k$ with the highest penalised
value. For candidates in value order and without a development bonus this is the exact optimum of the
node part. Flows are then fitted to the target shares and to demand by biproportional (RAS) scaling
of the edge matrix, $f_{pc}=r_pW_{pc}k_c$, which splits each consumer's demand in proportion to edge
quality and minimises the Kullback--Leibler divergence from the edge matrix among the matrices with
the required margins. Links below $\varepsilon$ of a consumer's demand are then pruned one at a time
and the fit repeated, but only where the consumer keeps $m_S$ suppliers and the refit still meets
the share targets: the share cap is a hard constraint, redundancy a soft one.

\begin{algorithm}[htbp]
\caption{Greedy fallback with water-filling (\textit{Greedy})}
\label{alg:greedy}
\begin{algorithmic}[1]
\Require nodes $C$, eligible set $\mathcal{E}$; $a,d,\delta$; complete edge matrix $W$; $\lambda>0$
\State $\bar W_p\gets\sum_cW_{pc}d_c$; \ $g_p\gets w_Na_p+w_E\bar W_p+w_D\delta_p$ ($a_p=0$ if undefined)
\State order $\mathcal{E}$ by $g_p$, $a_p$, $\bar W_p$ (decreasing) and code; move nodes that reach
  no consumer with demand to the end
\For{$k=m_P,\dots,|\mathcal{E}|$}
  \State $x^{(k)}\gets$ water-filling \eqref{eq:alg_waterfill} of the first $k$ nodes with
    $v_p=g_p$ (bisection on $\mu$)
  \State $V_k\gets\sum_pg_px^{(k)}_p+w_D\sum_p\delta_p-\lambda\sum_p(x^{(k)}_p)^2$
\EndFor
\State $k^\ast\gets\argmax_kV_k$; \ $S\gets$ the first $k^\ast$ nodes; \ $\sigma\gets x^{(k^\ast)}$
\State $f\gets$ RAS fit of $W_{S\times C}$ to row sums $\sigma$ and column sums $d$
  (at most 2{,}000 iterations, tolerance $10^{-12}$)
\State prune links with $f_{pc}<\varepsilon d_c$ one at a time where the consumer keeps $m_S$
  suppliers; refit; reject a refit that misses $\sigma$ or $d$ by more than $10^{-6}$
\State \Return $S$, $x=\sum_cf_{pc}$, $f$ and $J$ evaluated with eq.~\eqref{eq:net_obj}
\end{algorithmic}
\end{algorithm}

A greedy design that satisfies every constraint of the MILP is feasible for it, so an optimal MILP
solution scores at least as high. In the production configuration without caps the MILP objective
is at least the greedy one in all 22 stages, by 2.4\% in the median stage and 4.0\% at most, and the
greedy design is more concentrated (median index 0.154 against 0.139); with the caps on, solving
every stage greedily gives a median index of 0.160, above the calibration target, and a summed
objective 2.5\% lower. The fallback is a safeguard, not an equivalent; it is not triggered in the
production design.

% ---------------------------------------------------------------------------
\subsection{Portfolio caps}
\label{subsec:alg_portfolio}

Algorithm~\ref{alg:portfolio} enforces the caps of eq.~\eqref{eq:net_cap}. The caps are computed
once: GDP (latest finite WDI year, the coalition median if none) is transformed to
$\log(1+\mathrm{GDP})$, scaled to $[0,1]$ over the nodes, mapped linearly to $[2,10]$, rounded and
clamped, and the largest economy is set to ten. An appearance requires a positive share and a strict
rank of at most $K$, ties broken by producer score and then by code; with shares on the midpoints of
the tangent grid a dense rank would count every tied producer and let one stage contribute more than
$K$ appearances.

\begin{algorithm}[htbp]
\caption{Portfolio caps across stages (\textit{PortfolioCaps})}
\label{alg:portfolio}
\begin{algorithmic}[1]
\Require stage results $R$; exclusion sets $\mathcal{X}$; caps $\kappa$; $K=5$; $I_{\max}=30$
\For{$\iota=1,\dots,I_{\max}$}
  \State in every stage, $\rho_{i,ts}\gets$ strict rank of $x_{i,ts}$ (decreasing; ties by $a$, then code)
  \State $A_i\gets|\{(t,s):x_{i,ts}>0,\ \rho_{i,ts}\le K\}|$\Comment{eq.~\eqref{eq:net_topk}}
  \State $O\gets\{i:A_i>\kappa_i\}$; \textbf{if} $O=\emptyset$ \textbf{then} converged; \textbf{break}
  \ForAll{$i\in O$}\Comment{all offenders use the same snapshot of $R$}
    \State bar $i$ from the $A_i-\kappa_i$ counted stages with the smallest $x_{i,ts}\,a_{i,ts}$:
      $\mathcal{X}_{ts}\gets\mathcal{X}_{ts}\cup\{i\}$
  \EndFor
  \State re-solve every stage whose exclusion set changed\Comment{Algorithm~\ref{alg:stage}}
\EndFor
\State \Return $R$, $\mathcal{X}$, $(\kappa_i,A_i)_i$, convergence status, passes, violations left
\end{algorithmic}
\end{algorithm}

Exclusion sets only grow and a barred node remains a consumer, so at most $(1+I_{\max})$ times the
number of stages are solved. Within a pass the offenders' exclusions are chosen from one snapshot and
do not depend on the order in which offenders are visited; across passes the outcome is
path-dependent, because barring a node can promote another into a stage's top $K$. Compliance
requires $\sum_i\kappa_i\ge3S^\ast$ for $S^\ast$ solved stages, since every stage has at least three
producers within its top $K$; here $\sum_i\kappa_i=218\ge66$. In the production design the loop
converges after seven passes (62 stage solves) with no violation, and every run of the calibration
grid converges, in four to eight passes. A joint program with the linear constraints
$\sum_{(t,s)}y_{i,ts}\le\kappa_i$ would optimise the stages together, at the cost of a program 22
times as large; it is left for future work.

% ---------------------------------------------------------------------------
\subsection{Calibration of the diversification weight}
\label{subsec:alg_calibration}

\paragraph{Target.} For each stage with an HS6 basket, the nodes' observed export concentration is
$\HHI^{\mathrm{N}}_{ts}=\sum_{j\in\mathcal{V}}\big(X_{jts}/\sum_{k\in\mathcal{V}}X_{kts}\big)^2$,
where $X_{jts}$ are node $j$'s 2024 basket exports in the reconciled Comtrade matrix (importer
reports, else the exporter's mirror declaration $\times1.06$; unobserved cells count as zero). A
single-country node counts its exports to every other economy; the EU node counts only exports to
importers outside the EU, because trade among its members is internal to the node, as a country's
domestic sales are, and would otherwise make the EU the leading exporter of nearly every stage. The
target is the median over the 18 basket stages, 0.1556; the stage values range from 0.095
(Batteries Midstream) to 0.355.

\paragraph{Procedure.} For each $\lambda$ on the grid, the full production configuration is
re-solved: 37 nodes, 22 stages, caps on, MILP with the 300-second limit, up to 30 cap passes, and the
inputs assembled exactly as \texttt{scripts/30} assembles them. The rule selects the smallest grid
value whose median $\HHI^{\mathrm{D}}_{ts}$ over the 22 stages is at most the target. The run at the
selected value reproduces the stored production design to $5\times10^{-16}$ in every share, and the
production script warns if the designed median exceeds the observed one on its own inputs, so a
change of inputs that breaks the rule cannot pass unnoticed. Table~\ref{tab:alg_lambda_grid} lists
the 36 grid points.

\begin{table}[htbp]
\centering
\begin{threeparttable}
\caption{The calibration grid of the diversification weight}
\label{tab:alg_lambda_grid}
\footnotesize
\setlength{\tabcolsep}{3.5pt}
\begin{tabular}{@{}l r r r r r r r@{\hspace{14pt}}l r r r r r r r@{}}
\toprule
$\lambda$ & $\widetilde{\HHI}$ & Eff. & Max. & Roles & Floor & $\Delta$ & M &
$\lambda$ & $\widetilde{\HHI}$ & Eff. & Max. & Roles & Floor & $\Delta$ & M \\
\midrule
0    & .3450 & 2.90 & .400 &  79 &  6 &   0.00 & 22 & 0.86 & .1480 & 6.76 & .305 & 209 & 56 &  $-$6.39 & 22 \\
0.5  & .2156 & 4.64 & .400 & 149 & 35 &  $-$2.71 & 22 & 0.87 & .1493 & 6.70 & .305 & 209 & 54 &  $-$6.44 & 22 \\
0.70 & .1640 & 6.10 & .355 & 190 & 54 &  $-$5.02 & 22 & 0.88 & .1469 & 6.81 & .305 & 211 & 54 &  $-$6.54 & 22 \\
0.71 & .1630 & 6.14 & .345 & 192 & 54 &  $-$5.17 & 22 & 0.89 & .1414 & 7.07 & .302 & 212 & 54 &  $-$6.62 & 22 \\
0.72 & .1622 & 6.16 & .335 & 193 & 51 &  $-$5.27 & 22 & 0.90 & .1411 & 7.09 & .302 & 214 & 59 &  $-$6.70 & 22 \\
0.73 & .1630 & 6.14 & .335 & 194 & 53 &  $-$5.36 & 22 & 0.91 & .1423 & 7.03 & .295 & 213 & 55 &  $-$6.71 & 22 \\
0.74 & .1620 & 6.17 & .335 & 197 & 57 &  $-$5.50 & 22 & 0.92 & .1347 & 7.43 & .310 & 216 & 58 &  $-$6.86 & 21 \\
0.75 & .1593 & 6.28 & .335 & 197 & 58 &  $-$5.55 & 22 & 0.93 & .1323 & 7.56 & .309 & 217 & 61 &  $-$6.99 & 21 \\
0.76 & .1583 & 6.32 & .331 & 199 & 62 &  $-$5.66 & 22 & 0.94 & .1340 & 7.46 & .285 & 221 & 64 &  $-$7.07 & 22 \\
0.77 & .1558 & 6.42 & .325 & 200 & 61 &  $-$5.72 & 22 & 0.95 & .1331 & 7.51 & .305 & 221 & 62 &  $-$7.18 & 21 \\
\textbf{0.78} & \textbf{.1520} & 6.58 & .325 & 201 & 61 &  $-$5.80 & 22 & 0.96 & .1326 & 7.54 & .285 & 225 & 70 &  $-$7.28 & 22 \\
0.79 & .1485 & 6.74 & .325 & 203 & 63 &  $-$5.90 & 22 & 0.97 & .1343 & 7.45 & .282 & 225 & 66 &  $-$7.39 & 22 \\
0.80 & .1502 & 6.66 & .320 & 205 & 63 &  $-$6.03 & 22 & 0.98 & .1310 & 7.63 & .276 & 227 & 65 &  $-$7.44 & 22 \\
0.81 & .1500 & 6.66 & .315 & 205 & 54 &  $-$6.09 & 22 & 0.99 & .1288 & 7.77 & .300 & 226 & 65 &  $-$7.53 & 21 \\
0.82 & .1500 & 6.66 & .315 & 206 & 53 &  $-$6.14 & 22 & 1.00 & .1291 & 7.74 & .275 & 230 & 70 &  $-$7.58 & 22 \\
0.83 & .1492 & 6.70 & .315 & 205 & 55 &  $-$6.16 & 22 & 1.1  & .1196 & 8.36 & .257 & 238 & 63 &  $-$8.17 & 22 \\
0.84 & .1487 & 6.73 & .305 & 208 & 58 &  $-$6.26 & 22 & 1.2  & .1126 & 8.88 & .245 & 251 & 77 &  $-$8.96 & 22 \\
0.85 & .1501 & 6.66 & .305 & 207 & 56 &  $-$6.29 & 22 & 1.5  & .0945 & 10.58 & .215 & 280 & 92 & $-$10.64 & 22 \\
\addlinespace
\multicolumn{16}{@{}l}{\emph{Without portfolio caps:} $\lambda=0.78$: $\widetilde{\HHI}=.1393$, Eff.\ 7.18, Max.\ .312, Roles 214, Floor 57, $\Delta=-5.14$, M 22.}\\
\bottomrule
\end{tabular}
\begin{tablenotes}[flushleft]\footnotesize
\item \emph{Notes:} Each row re-solves the production configuration (37 nodes, 22 Upstream and
Midstream stages, portfolio caps on, MILP with a 300-second stage limit) at one value of $\lambda$.
$\widetilde{\HHI}$: median over the 22 stages of the Herfindahl index of designed shares; the rule
selects the smallest $\lambda$ with $\widetilde{\HHI}\le0.1556$, the median observed export
Herfindahl index of the nodes over the 18 basket stages (bold), and every row from 0.78 on meets
it. Eff.: median effective number of producers. Max.: largest designed share. Roles: producer roles
summed over stages. Floor: shares at $\underline s=0.05$. $\Delta$: summed node and edge terms of the
22 stage objectives (without the penalty), in percent of their value at $\lambda=0$ with caps. M:
stages solved by the MILP in the final pass; in the four runs with 21 (0.92, 0.93, 0.95, 0.99),
Nuclear Upstream reached the time limit and was solved by the greedy fallback. The
portfolio loop converged in every run (four to eight passes). Source: lambda study on the inputs of
the final rerun of 12 October 2026.
\end{tablenotes}
\end{threeparttable}
\end{table}

\paragraph{Results.} Without the penalty the median stage has three producers, two at the cap, and an
index of 0.345; at $\lambda=0.5$ it is still 0.216. Between 0.70 and 0.77 the median lies within
0.009 of the target and above it, and the rule selects 0.78 (0.1520, a margin of 0.0036). The median
rises in eight of the 35 grid steps (at 0.72--0.73, 0.79--0.80, 0.84--0.85, 0.86--0.87, 0.90--0.91,
0.93--0.94, 0.96--0.97 and 0.99--1.00), by at most 0.0017, and is flat between 0.81 and 0.82; no
point above the pick returns above the target. Every run up to $\lambda=0.91$ solves all 22 stages
by the MILP in every pass. The score cost of diversification grows smoothly, from 5.8\% at the pick
to 7.6\% at $\lambda=1$ and 10.6\% at 1.5, while the number of producer roles rises from 201 to 280.
Without caps, the same $\lambda$ gives a median of 0.139, so the calibrated value is specific to the
capped configuration, as it is to the edge weight $w_E=0.5$ (Section~\ref{sec:robustness}).

% ---------------------------------------------------------------------------
\subsection{Implementation, verification and options not used}
\label{subsec:alg_tests}

The stage MILP is written with \texttt{ompr} \citep{schumacher2023ompr} and passed through
\texttt{ompr.roi} and the ROI plug-in to GLPK \citep{makhorin2020glpk}. The production script writes
the tables \texttt{allied\_network\_*.csv} (specialization, flows, diagnostics, input checks, nodes,
edges, member shares, observed concentration, portfolio and build candidates); in the flow and edge tables
the reporter is always the producer and the partner the consumer, and the friendshore column holds
$\FS_{cp}$, the consumer's rating of the producer. The six test files of the design
(\texttt{tests/testthat/test-allied-network-*.R}), on synthetic fixtures only, run 43 tests with
443 passing expectations and no failure. They check the uniqueness and orientation of nodes and
edges, the sums of shares, demand and flows, the share bounds and redundancy rules, the tangent
approximation and its error bound, the water-filling fallback, the aggregation of the EU node and the
split of its shares, the strict top-$K$ count, the portfolio loop and the deterministic treatment of
ties.

\paragraph{Options not used in production.} Two options are retained in the code for other callers
and are not used for any result in the paper. The first is the three-component producer score with
the Policy index (node weights 0.5, 0.3 and 0.2 on $\EO$, policy and $\ES$), in which a node without
a policy value receives the stage median of the observed values before the S-curve,
\begin{equation}\label{eq:net_policy_impute}
  \POL_{pts}\ \leftarrow\ \operatorname{median}\{\POL_{qts}:\ q\in\mathcal{V},\ \POL_{qts}\
  \text{observed}\}\quad\text{if } \POL_{pts}\ \text{is missing};
\end{equation}
the production design switches the policy term off, so nothing is imputed. The second is a
capacity-denominated variant, in which flows are in US dollars against a stage target, producers have
existing and maximum capacities and capacity builds carry a cost that development potential can
lower. It needs stage targets and capacity data that do not exist at the resolution of the indices,
and its solvers lack the deterministic tie rule and the input checks of the share-based ones.
